% 来源：OEBPS/Text/Section0001_split_005.xhtml → image_2dcee32b8b11470b8378d41655b8bec8.jpg
\begin{longtable}{@{}C{\dimexpr 0.50000\linewidth-1.00000\tabcolsep\relax}C{\dimexpr 0.50000\linewidth-1.00000\tabcolsep\relax}@{}}
\caption{标准普尔指数对比全天候投资组合（1928—2013年）}\label{tab:27-2}\\
\toprule\noalign{\vskip 3pt}
{\heiti 标准普尔指数} & {\heiti 全天候投资组合} \\
\midrule
\endfirsthead
\multicolumn{2}{r}{续表}\\[0.6em]
\toprule\noalign{\vskip 3pt}
{\heiti 标准普尔指数} & {\heiti 全天候投资组合} \\
\midrule
\endhead
\bottomrule
\endfoot
\bottomrule
\endlastfoot
87年里标准普尔指数亏损年度24个（占比约27\%）\footnote{\label{fn:table-表27-2标准普尔指数对比全天候投资组合-1}包括红利再投资} & 同期全天候投资组合亏损年度只有14个（这意味着其余73年都是赚钱的） \\
20世纪30年代“大萧条”最严重的时候，1929—1932年标准普尔指数连续4年下跌，累计跌幅64.4\% & 同期全天候投资组合累计亏损20.55\%（优于标准普尔指数59\%） \\
平均年度亏损13.66\% & 平均年度亏损3.65\% \\
\end{longtable}
